\chapter{Lecture 20: Visitor Continued, Compilation Dependencies and Pimpl}
\begin{center}
\textit{March 20, 2018}
\end{center}

\begin{aside}
The original notes for this lecture are missing; this chapter reconstructs it
from the course order (Visitor in Lecture 19, exception safety in Lecture 21)
and the standard CS 246 sequence of that term.
\end{aside}

\section{Visitor Recap}
Last time: double dispatch.
\begin{lstlisting}[style=cpp]
e->strike(*w);
// 1. virtual Enemy::strike   -> chooses on the enemy's run-time type
// 2. w.useOn(*this)          -> overload chosen by the static type of *this
// 3. virtual Weapon::useOn   -> chooses on the weapon's run-time type
\end{lstlisting}

\section{Visitor, Second Use: Adding Operations to a Hierarchy}
The Visitor pattern can also be used to add functionality to an existing class
hierarchy \emph{without} adding new virtual methods to every class.

Suppose we want to do something with our Book hierarchy that its designers did
not anticipate --- e.g.\ build a catalogue: count Books by author, Texts by
topic, Comics by hero. Each new operation would normally mean a new virtual
method in \code{Book}, \code{Text} and \code{Comic}. Instead, add a single
\code{accept} method once, and let \emph{visitors} supply the operations:
\begin{lstlisting}[style=cpp]
class BookVisitor {
  public:
    virtual void visit(Book &b) = 0;
    virtual void visit(Text &t) = 0;
    virtual void visit(Comic &c) = 0;
    virtual ~BookVisitor() {}
};

class Book {
  public:
    virtual void accept(BookVisitor &v) { v.visit(*this); }
    virtual ~Book() {}
    ...
};
class Text : public Book {
  public:
    void accept(BookVisitor &v) override { v.visit(*this); }
    ...
};
class Comic : public Book {
  public:
    void accept(BookVisitor &v) override { v.visit(*this); }
    ...
};
\end{lstlisting}
(\code{accept} plays the role of \code{strike}; \code{visit} plays the role of
\code{useOn}.)

A visitor that builds the catalogue:
\begin{lstlisting}[style=cpp]
class Catalogue : public BookVisitor {
    std::map<std::string, int> theCatalogue;
  public:
    std::map<std::string, int> getCatalogue() const { return theCatalogue; }
    void visit(Book &b) override { ++theCatalogue[b.getAuthor()]; }
    void visit(Text &t) override { ++theCatalogue[t.getTopic()]; }
    void visit(Comic &c) override { ++theCatalogue[c.getHero()]; }
};

std::vector<Book*> collection{...};
Catalogue cat;
for (auto b : collection) b->accept(cat);
for (auto &p : cat.getCatalogue())
    std::cout << p.first << ": " << p.second << std::endl;
\end{lstlisting}
Recall that \code{m[key]} inserts \code{0} for a missing key, so
\code{++theCatalogue[...]} counts correctly from zero.

Any further operation (e.g.\ total pages per category, printing a report) is
just another \code{BookVisitor} subclass; \code{Book}, \code{Text} and
\code{Comic} are untouched.

\begin{important}[Trade-off of Visitor]
Visitor makes it easy to add new \emph{operations}, but hard to add new
\emph{classes} to the visited hierarchy: a new Book subclass requires a new
\code{visit} overload in \code{BookVisitor} and in every visitor.
\end{important}

\section{Compilation Dependencies}
Visitor creates a lot of mutual dependencies: \code{book.h} mentions
\code{BookVisitor}, and \code{bookvisitor.h} mentions \code{Book}, \code{Text}
and \code{Comic}. If each header \code{\#include}s the other we get an include
cycle, and the include guards cause one of the files to see an undefined class.

\paragraph{Q.} When do we need \code{\#include "a.h"}, and when does a
\textbf{forward declaration} \code{class A;} suffice?

The compiler needs the full definition of a class when it needs to know the
class's \emph{size} or its \emph{members}:
\begin{lstlisting}[style=cpp]
// a.h
class A { ... };

// b.h
#include "a.h"
class B : public A { ... };      // inheritance: needs A's size and members

// c.h
#include "a.h"
class C {
    A myA;                       // a field of type A: needs A's size
};

// d.h
class A;                         // forward declaration is enough
class D {
    A *myAp;                     // pointers all have the same size
};

// e.h
class A;
class E {
    A f(A x);                    // only a declaration: no size needed yet
};
\end{lstlisting}
In the implementation files, D and E will need the full definition of A
(\code{d.cc} dereferences the pointer; \code{e.cc} actually passes and returns
A's), so \code{d.cc} and \code{e.cc} \code{\#include "a.h"}.

\begin{important}[Advice]
Do not \code{\#include} a header when a forward declaration suffices. Prefer
forward declarations in headers and includes in \code{.cc} files.
\end{important}
Why bother?
\begin{itemize}
  \item it breaks include cycles;
  \item it reduces compilation dependencies: if \code{a.h} changes, only files
        that really include it are recompiled (\code{make} with \code{-MMD}
        follows the includes). Large projects can take hours to build, so this
        matters.
\end{itemize}

\section{The Pimpl Idiom}
Even a private field creates a compilation dependency:
\begin{lstlisting}[style=cpp]
// window.h
#include <X11/Xlib.h>
class XWindow {
    Display *d;
    Window w;
    int s;
    GC gc;
    unsigned long colours[10];
  public:
    ...
};
\end{lstlisting}
Every client that includes \code{window.h} also includes the whole X11 header,
and must be recompiled whenever a \emph{private} field of \code{XWindow} is added
or changed --- even though clients can't use those fields.

\begin{definition}
\textbf{Pimpl} (``pointer to implementation''): move all private fields into a
separate implementation struct, defined only in the \code{.cc} file, and store
only a pointer to it in the class.
\end{definition}

\begin{lstlisting}[style=cpp]
// window.h
struct XWindowImpl;              // forward declaration only
class XWindow {
    XWindowImpl *pImpl;
  public:
    XWindow();
    ~XWindow();
    ...                          // no change to the public interface
};

// windowimpl.h
#include <X11/Xlib.h>
struct XWindowImpl {
    Display *d;
    Window w;
    int s;
    GC gc;
    unsigned long colours[10];
};

// window.cc
#include "window.h"
#include "windowimpl.h"
XWindow::XWindow() : pImpl{new XWindowImpl} { ... }
XWindow::~XWindow() { delete pImpl; }
// other methods use pImpl->d, pImpl->w, ... instead of d, w, ...
\end{lstlisting}
Now clients include only \code{window.h}, which does not include X11. If the
implementation struct changes, only \code{window.cc} is recompiled; clients just
re-link. (Since the class now owns heap memory, the Rule of Five applies; with a
smart pointer from Lecture~21, \code{std::unique\_ptr<XWindowImpl>}, most of it
is automatic.)

\section{The Bridge Pattern}
Taking Pimpl further: if there are several possible window implementations
(X11, a Windows API, a test stub), make the implementation struct a
\emph{polymorphic class hierarchy}:
\begin{lstlisting}[style=cpp]
class WindowImpl {               // abstract implementation interface
  public:
    virtual void drawPixel(int x, int y, int colour) = 0;
    virtual ~WindowImpl() {}
};
class XWindowImpl : public WindowImpl { ... };
class MSWindowImpl : public WindowImpl { ... };

class Window {                   // abstraction
    WindowImpl *pImpl;           // chosen at construction time
  public:
    void fillRectangle(int x, int y, int w, int h, int colour);  // uses pImpl
};
\end{lstlisting}

\begin{center}
\begin{tikzpicture}[cls/.style={draw, minimum width=22mm, minimum height=8mm,
  font=\small}]
\node[cls] (w) {Window};
\node[cls, right=30mm of w] (wi) {\textit{WindowImpl}};
\node[cls, below left=12mm and -8mm of wi] (x) {XWindowImpl};
\node[cls, below right=12mm and -8mm of wi] (m) {MSWindowImpl};
\draw[-{Diamond[fill=black,length=4mm,width=2.5mm]}] (wi) -- (w);
\draw[-{Triangle[open,length=3mm,width=3mm]}] (x.north) -- (wi.south west);
\draw[-{Triangle[open,length=3mm,width=3mm]}] (m.north) -- (wi.south east);
\end{tikzpicture}
\end{center}

\begin{definition}
\textbf{Bridge pattern}: separate an abstraction from its implementation so the
two can vary independently --- new kinds of windows (subclassing
\code{Window}) and new platforms (subclassing \code{WindowImpl}) can each be
added without touching the other hierarchy.
\end{definition}

\section*{Summary}
\begin{itemize}
  \item Visitor also adds operations to a hierarchy: one \code{accept} per class,
        one visitor class per operation. Easy to add operations, hard to add
        classes.
  \item Include a header only when you need a class's size or members; otherwise
        forward-declare. This breaks cycles and speeds up rebuilds.
  \item Pimpl hides private fields behind a pointer to an implementation struct,
        so clients don't recompile when the implementation changes.
  \item Bridge: Pimpl with a polymorphic implementation hierarchy.
\end{itemize}
